% Generated draft: requires mathematical and source review.
\documentclass[11pt]{article}
\usepackage[margin=25mm]{geometry}
\usepackage{fontspec}
\setmainfont{Latin Modern Roman}
\usepackage{amsmath,amssymb,mathtools,unicode-math}
\setmathfont{Latin Modern Math}
\usepackage{enumitem,needspace,xurl,hyperref}
\setlength{\parindent}{0pt}
\setlength{\parskip}{5pt}
\setlength{\emergencystretch}{3em}
\newcommand{\N}{\mathbb{N}}
\newcommand{\Z}{\mathbb{Z}}
\newcommand{\Q}{\mathbb{Q}}
\newcommand{\R}{\mathbb{R}}
\newcommand{\C}{\mathbb{C}}
\newcommand{\F}{\mathbb{F}}
\newcommand{\A}{\mathbb{A}}
\newcommand{\PP}{\mathbb{P}}
\newcommand{\E}{\mathbb{E}}
\newcommand{\eps}{\varepsilon}
\newcommand{\dd}{\,\mathrm{d}}
\begin{document}
\section*{UnsolvedMath ID: 20000409}
\textbf{Source problem code:} catalogue entry
\subsection{Definitions and mathematical statement}
C. Hall asks for constructions of points on the Jacobian of a higher-genus hyperelliptic curve over $K=\F_q(T)$, analogous to the source's Drinfeld--Heegner construction for elliptic analytic rank one. For a smooth projective hyperelliptic curve $X/K$ of genus $g\ge2$, its Jacobian $J_X=\operatorname{Pic}^0(X)$ parametrizes degree-zero divisor classes. Analytic rank means the order of vanishing of its Hasse--Weil $L$-function at the central point. The question seeks rational classes expected from that analytic information; it does not assert a general equality of analytic and Mordell--Weil ranks or specify an algorithm and independence guarantee. The record's comment is truncated after ``when the analytic rank is,'' and cannot be completed from the supplied verification.

The attached unreviewed aid treats only a split genus-two family. Assume $K$ has odd characteristic and $F(u)\in K[u]$ is a separable cubic with $F(0)\ne0$. Let $X,E_1,E_2$ be the smooth projective models of
\[
X:y^2=F(x^2),\qquad E_1:v^2=F(u),\qquad E_2:w^2=uF(u).
\]
The first curve has genus two. The elliptic models have the usual point at infinity $O_1$ on $E_1$ and a rational point $(0,0)$ on the quartic model of $E_2$, allowing an elliptic origin. The degree-two quotient maps are
\[
\pi_1(x,y)=(x^2,y),\qquad\pi_2(x,y)=(x^2,xy).
\]
The aid claims these induce a $K$-isogeny $J_X\to E_1\times E_2$.

Given an already constructed non-torsion Drinfeld--Heegner input $P\in E_1(K)$, it defines
\[
D_P=\pi_1^*([P-O_1])\in J_X(K),\qquad
\pi_{1,*}D_P=2[P-O_1].
\]
The pullback is scheme-theoretic: individual points above $P$ need not be $K$-rational, but the pulled-back divisor class is. The displayed norm relation certifies non-torsion. The aid reports $L(J_X,s)=L(E_1,s)L(E_2,s)$, so analytic ranks one and zero for $E_1,E_2$ give analytic rank one for $J_X$. Existence and nonvanishing of the required Heegner input need their own hypotheses. This transfer construction is confined to the split genus-two locus and does not solve the general higher-genus construction problem.
\subsection{Short English statement}
Construct rational Jacobian points suggested by analytic rank, beyond elliptic curves. The aid transfers an assumed non-torsion Heegner point through an explicit degree-two quotient on a special split genus-two family.
\subsection{Sources}
\begin{itemize}
\item C. Hall, Problem 1, \emph{Problems from the AIM Tate Conjecture Workshop}, July 23--27, 2007, transcribed by Christopher Lyons: \url{https://aimath.org/WWN/tateconjecture/tateconjecture.pdf}. Canonical location: \texttt{aim-algebraic-number-theory-notes.json}, notes[60]. The supplied record is reconstructed/unverified and truncated; its verification announces a complete source record without reproducing it. No PDF was consulted.
\item The original invokes Drinfeld Heegner points without a full reference. The attached split-Jacobian aid is machine generated and not peer reviewed, with low novelty confidence and no complete authoritative bibliography.
\item Local registry: \texttt{data/unsolvedmath/problems.json}, permanent catalogue entry, exact source code \texttt{catalogue entry}.
\end{itemize}
\end{document}
